\chapter{Cataract Surgery}

\section*{What Is This Surgery?}
Cataract surgery is a prevalent and highly effective ophthalmic procedure aimed at restoring vision by removing the eye's natural crystalline lens when it has become opacified due to cataracts. This condition, characterized by clouding of the lens, leads to significant visual impairment, affecting daily activities such as reading, driving, and recognizing faces. The surgical intervention involves the meticulous extraction of the opaque lens and its replacement with a clear, artificial intraocular lens (IOL). The primary anatomical site of this surgery is the crystalline lens, located behind the iris and pupil within the posterior chamber of the eye. Given its high success rate and the profound impact on patients' quality of life, cataract surgery is one of the most commonly performed surgical procedures worldwide.

\section*{Alternative Names}
\begin{itemize}
  \item Lens Replacement Surgery
  \item Phacoemulsification (Phaco)
  \item Extracapsular Cataract Extraction (ECCE)
  \item Intracapsular Cataract Extraction (ICCE) -- largely historical
  \item Intraocular Lens Implantation
  \item Clear Lens Exchange (when performed for refractive purposes without significant cataract)
\end{itemize}

\section*{Relevant Surgical Anatomy}
The crystalline lens is a transparent, biconvex structure located in the anterior segment of the eye, posterior to the iris and anterior to the vitreous body. It is suspended by delicate zonular fibers that originate from the ciliary body and attach to the lens capsule, allowing for accommodation, which is the ability of the lens to change shape for focusing. The lens is encased within a thin, elastic capsule, comprising an anterior and a posterior layer. During cataract surgery, particularly phacoemulsification, the anterior capsule is opened (capsulorhexis), and the lens material is removed while leaving the posterior capsule intact to support the implanted IOL.

Key anatomical structures relevant to cataract surgery include:

\begin{itemize}
  \item \textbf{Cornea:} The clear, outer layer of the eye through which incisions are made.
  \item \textbf{Anterior Chamber:} The space between the cornea and the iris, filled with aqueous humor.
  \item \textbf{Iris:} The colored part of the eye that controls pupil size and must be carefully managed during surgery.
  \item \textbf{Ciliary Body:} Responsible for aqueous humor production and lens accommodation, located adjacent to the lens.
  \item \textbf{Zonules:} Fibers that connect the ciliary body to the lens, crucial for lens stability.
\end{itemize}

The integrity of the zonules is paramount for stable IOL placement, and the rich vascular and nerve supply of the surrounding ciliary body and iris are important considerations for anesthesia and managing intraoperative complications.

\section*{Surgical Principle \& Rationale}
The fundamental principle of cataract surgery is to restore clarity to the visual axis by excising the opacified natural lens and replacing it with a transparent, precisely calculated artificial lens. The most common technique employed today is phacoemulsification, which utilizes high-frequency ultrasound energy to emulsify the hard nucleus of the cataract into small fragments. These fragments are then aspirated from the eye through a small incision. The posterior capsule of the lens is typically left intact, serving as a stable platform for the implantation of a folded IOL.

The rationale behind this surgical intervention is to correct the optical defect caused by the cataract, which scatters light and reduces visual acuity. By replacing the cloudy lens with a clear IOL, light can be focused sharply onto the retina, restoring clear vision. Technical goals include achieving a clear visual axis, ensuring stable and centered IOL placement within the capsular bag, minimizing surgically induced astigmatism, and optimizing the patient's postoperative refractive outcome. Modern IOL technology offers various options, including monofocal lenses, toric lenses for astigmatism correction, and multifocal or extended depth of focus (EDOF) lenses, allowing for customized visual rehabilitation.

\section*{History \& Surgical Evolution}
The history of cataract surgery dates back to antiquity, with the earliest known attempts to treat cataracts through "couching" practiced in ancient India and Egypt around 800 BC. This method involved dislodging the cloudy lens into the vitreous cavity using a sharp instrument, providing temporary visual improvement but often leading to severe complications such as infection and blindness. A significant milestone occurred in 1747 when French ophthalmologist Jacques Daviel performed the first recorded successful extracapsular cataract extraction (ECCE), where the lens was removed through an incision while leaving the posterior capsule intact.

The introduction of the intraocular lens (IOL) by Sir Harold Ridley in 1949 marked another pivotal moment in cataract surgery. Ridley observed that Royal Air Force pilots with acrylic fragments in their eyes during World War II did not experience significant inflammatory reactions, leading to the development of the first IOL. The next major advancement came in 1967 with Charles Kelman's introduction of phacoemulsification, a technique that uses ultrasound to break up the lens, allowing its removal through a much smaller incision. This innovation dramatically reduced recovery times and improved outcomes, paving the way for the micro-incisional, highly refined procedures performed today, often incorporating femtosecond laser assistance for enhanced precision.

\section*{Clinical Indications}
Cataract surgery is indicated when the opacification of the crystalline lens leads to visual impairment that significantly impacts a patient's quality of life or poses other ocular health risks. The following clinical indications warrant surgical intervention:

\begin{itemize}
  \item Significant reduction in best-corrected visual acuity (e.g., 20/40 or worse) that cannot be improved with glasses or contact lenses.
  \item Presence of debilitating glare, halos, or starbursts around lights, particularly at night, which impairs driving or other activities.
  \item Difficulty performing daily activities such as reading, driving, watching television, or recognizing faces due to blurred or dim vision.
  \item Impaired vision that affects occupational performance, hobbies, or overall functional independence.
  \item Secondary glaucoma, such as phacomorphic glaucoma (due to lens swelling) or phacolytic glaucoma (from leaking lens proteins), where lens removal is therapeutic.
  \item Lens-induced uveitis, an inflammatory condition caused by leakage of lens material into the anterior chamber.
  \item When a clear view of the retina is required for the diagnosis or treatment of other posterior segment diseases, such as diabetic retinopathy, macular degeneration, or retinal detachment.
  \item Significant anisometropia (a large difference in refractive error between the two eyes) caused by a cataract, which cannot be corrected with spectacles.
  \item Monocular diplopia (double vision in one eye) caused by irregular astigmatism or opacities within the cataract.
  \item In some cases, for purely refractive purposes (clear lens exchange) in patients who desire to reduce dependence on glasses, particularly for presbyopia correction, though this is not a primary indication for cataract removal.
\end{itemize}

\section*{Clinical Positioning}
Cataract surgery is overwhelmingly considered the primary and definitive curative treatment for symptomatic cataracts. Once a cataract develops to the point of causing visual impairment, there are no medications, eye drops, or lifestyle changes that can reverse its progression or restore lens clarity. Therefore, surgical removal and IOL implantation are the only effective solutions, making it a first-line, primary intervention.

In cases where a cataract coexists with other significant ocular pathologies, such as advanced glaucoma, corneal disease, or retinal detachment, cataract surgery may be performed as part of a broader, multi-stage treatment plan. While still addressing the cataract primarily, its timing and specific technique might be influenced by the need to facilitate or complement other necessary ocular interventions, such as combined cataract and glaucoma surgery (phaco-trabeculectomy) or vitrectomy for retinal conditions. In these scenarios, it serves as a crucial component of a comprehensive approach to restore overall ocular health and function.

\section*{Surgical Technique \& Procedure}
Cataract surgery is typically performed on an outpatient basis, usually under topical anesthetic drops combined with a mild oral or intravenous sedative to ensure patient comfort and cooperation. In some cases, a local anesthetic injection (peribulbar or retrobulbar block) may be used for more profound anesthesia and akinesia (paralysis of eye movement). The patient is positioned supine, and the eye and surrounding area are prepped with an antiseptic solution (e.g., povidone-iodine) and draped to maintain a sterile field. A lid speculum is used to keep the eyelids open.

The procedure generally follows these steps:

\begin{enumerate}
  \item \textbf{Incision Creation:} One or two small, self-sealing incisions (typically 1.8-3.0 mm) are made at the limbus (the junction of the cornea and sclera) or in the clear cornea. These incisions are designed to be watertight without sutures.
  \item \textbf{Viscoelastic Injection:} A viscoelastic substance, a clear, gel-like material, is injected into the anterior chamber to maintain space, protect delicate intraocular structures (like the corneal endothelium), and facilitate subsequent steps.
  \item \textbf{Capsulorhexis:} A precise, continuous curvilinear opening is created in the anterior capsule of the lens using a cystotome or forceps. This opening, known as a continuous curvilinear capsulorhexis (CCC), is crucial for stable and centered IOL placement within the capsular bag.
  \item \textbf{Hydrodissection and Hydrodelineation:} Balanced salt solution is injected between the lens capsule and the cortex (hydrodissection), and between the nucleus and epinucleus (hydrodelineation), to separate these layers and allow for easier rotation and removal of the lens material.
  \item \textbf{Phacoemulsification:} An ultrasound probe is inserted through the main incision. The probe's tip vibrates at ultrasonic frequencies, emulsifying the hard lens nucleus into small fragments. The fragmented material is then aspirated from the eye through the same probe.
  \item \textbf{Cortical Aspiration:} The softer cortical material remaining after phacoemulsification is removed from the capsular bag using an irrigation/aspiration (I/A) handpiece, ensuring a clean and clear capsular bag.
  \item \textbf{IOL Insertion:} A folded intraocular lens, chosen based on precise preoperative biometry measurements, is injected through the small incision using a specialized injector. Once inside the eye, the IOL unfolds and is carefully positioned within the capsular bag, ensuring proper centration and orientation.
  \item \textbf{Viscoelastic Removal:} The remaining viscoelastic material is thoroughly aspirated from the eye to prevent postoperative intraocular pressure spikes.
  \item \textbf{Wound Hydration/Closure:} The corneal incisions are typically self-sealing due to their architectural design. The surgeon may hydrate the corneal stroma around the incision to ensure a watertight seal. Sutures are rarely required for modern micro-incisional phacoemulsification.
\end{enumerate}

At the end of the procedure, antibiotic and anti-inflammatory drops are often applied, and a protective eye shield may be placed to prevent accidental rubbing or trauma.

\section*{Preoperative Workup \& Preparation}
Preparation for cataract surgery involves several key steps to ensure patient safety and optimal outcomes. A comprehensive preoperative ophthalmologic examination is performed, including detailed visual acuity testing, slit-lamp examination of the anterior segment, fundoscopy to assess retinal health, and precise biometry measurements to calculate the correct power of the intraocular lens. Systemic medical evaluation may be required, especially for patients with significant comorbidities (e.g., cardiac disease, diabetes), to obtain medical and anesthesia clearance, often involving a primary care physician or internist.

Specific preparatory measures include:

\begin{itemize}
  \item \textbf{Medication Review:} Patients are instructed to continue most systemic medications, but adjustments may be needed for blood thinners (often continued, but discussed with the surgeon and cardiologist) or diabetes medications. Alpha-blockers (e.g., tamsulosin) must be disclosed due to the risk of Intraoperative Floppy Iris Syndrome (IFIS).
  \item \textbf{Preoperative Eye Drops:} Patients typically begin a regimen of antibiotic and anti-inflammatory eye drops a few days prior to surgery to reduce the risk of infection and inflammation.
  \item \textbf{NPO Status:} Fasting from food and drink for at least 6-8 hours before surgery is required if sedation or general anesthesia is planned.
  \item \textbf{Informed Consent:} Detailed discussion of the procedure, potential risks, benefits, alternative treatments, and IOL options, followed by signing an informed consent form.
  \item \textbf{Transportation:} Arranging for someone to drive the patient home after the procedure, as vision may be blurry and sedation effects may linger.
  \item \textbf{Prophylactic Antibiotics:} Intracameral antibiotics (e.g., cefuroxime) are often administered at the end of surgery to further reduce the risk of endophthalmitis.
\end{itemize}

\section*{Contraindications \& Precautions}
While cataract surgery is generally safe, certain conditions may contraindicate or necessitate precautions for the procedure.

\textbf{Absolute Contraindications:}

\begin{itemize}
  \item Active ocular infection (e.g., endophthalmitis, severe conjunctivitis, keratitis) that has not been adequately treated, as surgery could exacerbate the infection.
  \item Uncontrolled severe systemic disease (e.g., unstable angina, recent myocardial infarction, severe respiratory distress, uncontrolled hypertension) that makes lying flat or tolerating anesthesia unsafe.
  \item Patient's inability to provide informed consent or cooperate during the procedure, which is crucial for successful outcomes under local anesthesia.
  \item No light perception in the affected eye, indicating no potential for visual improvement.
  \item Unrealistic patient expectations that cannot be met by the surgery.
\end{itemize}

\textbf{Relative Contraindications \& Precautions:}

\begin{itemize}
  \item Uncontrolled glaucoma or ocular hypertension, which should ideally be managed before surgery to minimize postoperative pressure spikes and further optic nerve damage.
  \item Severe dry eye syndrome or significant blepharitis, which should be treated preoperatively to reduce infection risk, improve comfort, and optimize corneal surface for biometry.
  \item Corneal pathology (e.g., Fuchs' dystrophy, significant corneal scarring, severe dry eye) that may compromise surgical visualization, increase the risk of corneal decompensation, or affect postoperative corneal clarity.
  \item Retinal pathology (e.g., active retinal detachment, severe macular edema, proliferative diabetic retinopathy) that may limit visual recovery despite successful cataract removal, requiring careful patient counseling.
  \item Use of alpha-blockers (e.g., tamsulosin for benign prostatic hyperplasia), which can cause Intraoperative Floppy Iris Syndrome (IFIS), requiring specific surgical techniques (e.g., iris hooks, pupil expansion rings) and precautions.
  \item Monocular patient, where the risks of surgery are weighed more heavily due to reliance on the single functioning eye.
  \item Pregnancy, where elective surgery is generally deferred until after delivery unless medically urgent due to the risks of anesthesia and medications.
  \item History of previous ocular trauma or surgery, which may complicate the procedure due to altered anatomy or zonular weakness.
  \item Extremely dense or "white" cataracts, which may increase surgical complexity and risk of complications.
\end{itemize}

\section*{Complications \& Risks}
While cataract surgery is generally safe and highly successful, like any surgical procedure, it carries potential risks. These can be broadly categorized as general surgical complications, anesthesia-related risks, and procedure-specific complications.

\begin{itemize}
  \item \textbf{General Surgical Risks:}
    \begin{itemize}
      \item \textbf{Infection (Endophthalmitis):} A rare but severe intraocular infection, occurring in less than 0.1\% of cases, which can lead to permanent vision loss if not promptly treated.
      \item \textbf{Bleeding:} May manifest as a hyphema (blood in the anterior chamber) or, more rarely, a retrobulbar hemorrhage if a local anesthetic injection is used, potentially causing orbital compression.
      \item \textbf{Inflammation:} Postoperative uveitis or iritis, usually managed with topical steroids, but can be severe.
    \end{itemize}
  \item \textbf{Anesthesia-Related Risks:}
    \begin{itemize}
      \item Allergic reactions to medications (e.g., local anesthetics, sedatives).
      \item Respiratory depression, cardiovascular events (e.g., arrhythmia, hypotension), or nausea/vomiting, particularly with sedation or general anesthesia.
      \item Ocular perforation or optic nerve damage from local anesthetic injections (extremely rare).
    \end{itemize}
  \item \textbf{Procedure-Specific Risks:}
    \begin{itemize}
      \item \textbf{Posterior Capsule Rupture (PCR) with Vitreous Loss:} The posterior capsule, which supports the IOL, can tear, potentially leading to vitreous prolapse into the anterior chamber, IOL malposition, or retinal complications such as detachment.
      \item \textbf{Retained Lens Fragments:} Small pieces of the lens nucleus or cortex may remain in the eye, potentially causing inflammation, elevated intraocular pressure, or requiring further surgical intervention.
      \item \textbf{Retinal Detachment:} A rare but serious complication where the retina separates from its underlying support tissue, requiring further surgery (vitrectomy).
      \item \textbf{Cystoid Macular Edema (CME):} Swelling of the macula (the central part of the retina), which can cause blurred vision and may require medical treatment with topical or injectable anti-inflammatory agents.
      \item \textbf{Corneal Edema/Decompensation:} Swelling of the cornea, which can be temporary due to surgical trauma or, in rare cases, permanent, requiring a corneal transplant (e.g., Descemet's membrane endothelial keratoplasty - DMEK).
      \item \textbf{Intraocular Lens (IOL) Dislocation or Decentration:} The artificial lens may shift from its intended position within the capsular bag, requiring repositioning or exchange.
      \item \textbf{Posterior Capsule Opacification (PCO):} Also known as "secondary cataract," this is the most common long-term complication, where residual lens epithelial cells proliferate on the posterior capsule, causing clouding and blurred vision. It is effectively treated with a quick, non-invasive YAG laser capsulotomy.
      \item \textbf{Refractive Error:} Despite careful biometry, residual astigmatism, myopia, or hyperopia may persist, requiring glasses or contact lenses for optimal vision.
      \item \textbf{Increased Intraocular Pressure (IOP):} Temporary or persistent elevation of eye pressure, which may require medication or, rarely, further surgery.
      \item \textbf{Diplopia (Double Vision):} Can occur if the IOL is decentered, if pre-existing ocular muscle imbalances are unmasked, or due to aniseikonia.
      \item \textbf{Ptosis:} Drooping of the eyelid, usually temporary due to speculum pressure or anesthetic effects, but can be persistent.
      \item \textbf{Dysphotopsia:} Perception of unwanted optical phenomena, such as glare, halos, or shadows, particularly with certain IOL designs.
    \end{itemize}
\end{itemize}

\section*{Postoperative Recovery Timeline}
Immediately after cataract surgery, the patient is moved to a recovery area (PACU) where vital signs are monitored. The operated eye may be covered with a protective shield, and initial instructions regarding eye drops and activity restrictions are provided. Vision in the operated eye will likely be blurry due to dilation, residual anesthetic, and the eye's initial adjustment to the new lens. Mild discomfort, grittiness, or a foreign body sensation is common. Patients are typically discharged within a few hours, provided they are stable and have arranged for transportation.

Over the first 24-48 hours, patients are advised to rest and avoid strenuous activities. They will begin a regimen of prescribed antibiotic and anti-inflammatory eye drops, which are crucial for preventing infection and controlling inflammation. Vision gradually improves over the first few days to a week, with colors often appearing brighter and sharper. Patients are instructed to avoid rubbing the eye, bending over, lifting heavy objects (typically >10-15 lbs), or swimming for at least one to two weeks to prevent complications. Full visual stabilization, including the final refractive outcome, typically occurs within one month, at which point a new glasses prescription, if needed, can be finalized.

\section*{Surgical Findings \& Pathology}
During cataract surgery, the surgeon documents several key findings. These include the type and density of the cataract (e.g., nuclear sclerotic, cortical, posterior subcapsular), the integrity of the zonular fibers, the ease of capsulorhexis and lens removal, and the presence of any intraoperative complications such as posterior capsule rupture or vitreous loss. The type, power, and specific placement of the intraocular lens (IOL) are also meticulously recorded.

In uncomplicated cataract surgery, the removed crystalline lens material is typically not sent for routine histopathological examination. The diagnosis of cataract is made clinically, and the lens material itself rarely harbors pathology requiring microscopic analysis. However, in specific circumstances, such as suspected phacomorphic glaucoma with an unusually swollen lens, or if there is a rare suspicion of an intraocular tumor involving the lens, the excised lens material may be sent to pathology for further evaluation. For the vast majority of cases, the "pathology" report is simply a confirmation of the removed lens material, with the primary focus of the surgical findings being the intraoperative observations and the successful implantation of the IOL.

\section*{Expected Postoperative Course}
A normal and successful postoperative course following cataract surgery is characterized by a progressive improvement in visual function and a smooth, uneventful healing process. Patients typically experience a significant reduction in glare, halos, and blurriness within days to weeks after the procedure. The eye should appear quiet, with minimal redness or inflammation, and the cornea should remain clear. Discomfort, grittiness, or a foreign body sensation usually subsides within the first few days. The intraocular lens should be well-centered within the capsular bag, and intraocular pressure should be within normal limits. Most patients achieve their target refractive outcome, often leading to significantly improved uncorrected or best-corrected visual acuity, enhancing their ability to perform daily activities and improving overall quality of life. The full visual potential is usually realized within 4-6 weeks as the eye completely heals and adjusts to the new lens.

\section*{Postoperative Care \& Follow-up}
Postoperative care is crucial for monitoring healing, preventing complications, and optimizing visual outcomes.

\begin{itemize}
  \item \textbf{Ophthalmology Visits:} Typically scheduled for post-op day 1, week 1, and approximately 4-6 weeks after surgery. These visits assess visual acuity, intraocular pressure, corneal clarity, IOL position, and overall ocular health. Additional visits may be scheduled if complications arise.
  \item \textbf{Eye Drop Regimen:} Patients continue prescribed antibiotic and anti-inflammatory eye drops for several weeks, with a gradual tapering schedule as directed by the surgeon. Strict adherence to this regimen is vital.
  \item \textbf{Activity Resumption:} Gradual return to normal activities is advised, with specific restrictions on heavy lifting, bending, rubbing the eye, and swimming for the initial weeks to prevent increased intraocular pressure or infection.
  \item \textbf{Glasses Prescription:} A final glasses prescription, if needed for optimal distance, near, or intermediate vision, is typically provided around 4-6 weeks post-surgery once the eye has fully stabilized and the refractive outcome is final.
  \item \textbf{Warning Signs:} Patients are educated on warning signs of potential complications, such as sudden vision loss, severe pain, increasing redness, discharge, new floaters, or flashes of light, and instructed to contact their surgeon immediately if these occur.
\end{itemize}

Regular follow-up ensures any long-term complications, such as posterior capsule opacification, are detected and managed promptly.

\section*{Epidemiology}
Cataracts are the leading cause of blindness and visual impairment worldwide, accounting for approximately 51\% of global blindness and a significant proportion of moderate to severe visual impairment. The incidence of cataracts increases significantly with age, making it predominantly a disease of the elderly. In developed countries, age-related cataracts are the most common indication for surgery, with the prevalence rising sharply after the age of 60, affecting over half of individuals by age 75. While both sexes are affected, some epidemiological studies suggest a slightly higher prevalence in women. Globally, over 30 million cataract surgeries are performed annually, making it one of the most common surgical procedures. The overall mortality rate associated with cataract surgery is extremely low, primarily linked to systemic comorbidities in elderly patients undergoing sedation or general anesthesia. Serious ocular morbidity, such as endophthalmitis, is rare, occurring in less than 0.1\% of cases, reflecting the high safety profile of modern techniques.

\section*{Outcomes \& Prognosis}
The outcomes of modern cataract surgery are overwhelmingly positive, with typical success rates exceeding 95\% in terms of achieving improved visual acuity. Most patients experience a significant enhancement in their functional vision, leading to improved quality of life, greater independence, and a reduced risk of falls. The prognosis for visual recovery is excellent, particularly in eyes without pre-existing ocular pathologies such as advanced macular degeneration or severe glaucoma. While the cataract itself cannot recur, the most common long-term complication is posterior capsule opacification (PCO), affecting 20-50\% of patients within five years. This "secondary cataract" is effectively treated with a quick, non-invasive YAG laser capsulotomy, which restores clear vision. Prognostic factors for optimal outcomes include the absence of other significant ocular diseases, accurate preoperative biometry and IOL power calculation, meticulous surgical technique, and patient compliance with postoperative medication and instructions. Studies consistently demonstrate the cost-effectiveness and high patient satisfaction associated with cataract surgery.

\section*{References}
\begin{enumerate}
  \item MedlinePlus. Cataract Surgery. U.S. National Library of Medicine; 2024. Available from: \url{https://medlineplus.gov/cataractsurgery.html}
  \item Kanski JJ, Bowling B. Kanski's Clinical Ophthalmology: A Systematic Approach. 9th ed. Elsevier; 2020.
  \item American Academy of Ophthalmology. Cataract in the Adult Eye Preferred Practice Pattern. San Francisco, CA: American Academy of Ophthalmology; 2021.
  \item Yanoff M, Duker JS. Ophthalmology. 6th ed. Elsevier; 2023.
\end{enumerate}

\clearpage